\documentclass[11pt,a4paper]{article}
\usepackage[T1]{fontenc}
\usepackage[utf8]{inputenc}
\usepackage[spanish]{babel}
\usepackage{amsmath,amssymb}
\usepackage[margin=2.5cm]{geometry}
\usepackage[hidelinks]{hyperref}
\newcommand{\B}{\{0,1\}}
\title{Murdoku: modelos del caso de uso\\Generar un puzzle}
\author{Documentaci\'on de la implementaci\'on}
\date{}

\begin{document}
\maketitle

\section{Objetivo y flujo}

El caso de uso \texttt{generate-puzzle}, ejecutado mediante
\texttt{murdoku-generate}, construye una soluci\'on objetivo $S^\star$,
enumera un pool de pistas verdaderas y selecciona un subconjunto que:
\begin{enumerate}
  \item garantice que $S^\star$ sea la \'unica colocaci\'on posible;
  \item minimice la distancia a la dificultad solicitada;
  \item a igualdad de distancia, use el menor n\'umero de pistas.
\end{enumerate}
La optimizaci\'on es condicional al mapa, las personas y el pool de
esta ejecuci\'on. No compara mapas distintos ni vuelve a sortear
$S^\star$ si el pool resulta insuficiente.

Se construyen tres clases de modelos:
\[
 \text{generaci\'on espacial}
 \ \longrightarrow\ \text{maestro de selecci\'on}
 \ \longleftrightarrow\ \text{subproblema de unicidad}.
\]
El pool se construye por enumeraci\'on y evaluaci\'on, no mediante
otro modelo CP-SAT.

Las fuentes principales son:
\begin{itemize}
  \item \path{murdoku_checker\mains\main_generate_new_puzzle.py}.
  \item \path{murdoku_checker\apps\solution_creator\placement_model.py}.
  \item \path{murdoku_checker\apps\pool_generator\candidate_rules.py}
        y \path{murdoku_checker\apps\pool_generator\pool_generator.py}.
  \item \path{murdoku_checker\apps\problem_builder\problem_builder.py}.
  \item \path{murdoku_checker\apps\solution_checker\solution_checker.py}.
\end{itemize}

\section{Fase 1: modelo espacial para obtener $S^\star$}

Sean $H,W$ las dimensiones, $N=HW$, $K$ las habitaciones y
$P=\{1,\ldots,n\}$ las personas, con $n=\min(H,W)$. La v\'ictima
es $v=n$. La especificaci\'on aleatoria fija tipos de muebles,
tama\~nos $s_t\in\{1,2\}$, cantidades $m_t\geq1$ y ocupabilidad $o_t$,
con presupuesto
\[
 \sum_t s_tm_t
   =\operatorname{round}\left(HW\,\frac{\alpha}{100}\right).
\]
Estos valores se fijan antes de CP-SAT.

El modelo espacial contiene:
\begin{description}
  \item[Habitaciones.] Una variable entera $r_c\in\{0,\ldots,K-1\}$
        por celda, una semilla fija $s_k$ por habitaci\'on,
        una profundidad $d_c\in\{0,\ldots,N-1\}$ y arcos
        booleanos $a_{cu}$ a vecinos ortogonales.
  \item[Muebles.] Un booleano $f_{iA}$ por ejemplar $i$ y huella
        admisible $A$, y un indicador $b_c$ de celda bloqueada.
  \item[Personas.] Fila $Y_p$, columna $X_p$, \'indice
        $Q_p=WY_p+X_p$, habitaci\'on $K_p$ y un booleano $h_p$
        que indica si $p\neq v$ acompa\~na a la v\'ictima.
\end{description}

Con $V$ el conjunto de celdas y $\mathcal N(c)$ sus vecinos,
se eligen semillas separadas antes de resolver. Sea
$S=\{s_k:k\in\mathcal K\}$. Las restricciones de conectividad son
\begin{align*}
 r_{s_k}&=k,\quad d_{s_k}=0,\\
 \sum_{u\in\mathcal N(c)}a_{cu}&=1\quad(c\notin S),&
 a_{cu}=1&\Longrightarrow r_c=r_u,\ d_c>d_u.
\end{align*}
Imponen habitaciones no vac\'ias y conexas, sin un m\'inimo de
cuatro celdas.
Las formas pueden ser no rectangulares: cada intersecci\'on con
una fila o columna es continua y el \'area ocupa al menos el 75\%
de su envolvente, con proporci\'on entre ejes como m\'aximo 1.5.
Se minimiza la distancia a las semillas con costes aleatorios
fijos, aceptando soluciones factibles tras un m\'aximo de diez
segundos sin exigir optimalidad espacial.

Adem\'as, para cada habitaci\'on se obtiene un indicador exacto
\[
 t_k=
 \left[
 \left|\{i(c):r_c=k\}\right|=1
 \ \lor\
 \left|\{j(c):r_c=k\}\right|=1
 \right],
 \qquad \sum_k t_k\leq1.
\]
Por tanto, como m\'aximo una habitaci\'on puede caber en una sola
fila o columna, contando las orientaciones juntas. Una celda
individual cuenta una vez y no se exige que haya ninguna habitaci\'on
estrecha. Se calculan adem\'as el n\'umero de filas $H_k$,
columnas $W_k$ y celdas $A_k$ de cada habitaci\'on, y se exige
\[
 4A_k\geq3H_kW_k,\qquad 2H_k\leq3W_k,\qquad 2W_k\leq3H_k.
\]
La conectividad hace que las coordenadas ocupadas sean consecutivas.
La desigualdad de \'area exige ocupar al menos el 75\% del
rect\'angulo envolvente y las otras limitan su proporci\'on a $3/2$.
Se permiten formas en L, pero no huecos. Una habitaci\'on estrecha
solo puede ser $1\times1$, y sigue siendo opcional y \'unica como
m\'aximo. La formulaci\'on auxiliar completa est\'a en
\texttt{generate-solution.tex}. Algunas configuraciones pueden
resultar infactibles bajo estas restricciones de forma.

Cada ejemplar elige una huella, no hay solapamientos, los muebles de
dos celdas quedan en una misma habitaci\'on y $b_c$ es la suma de
las huellas seleccionadas no ocupables que cubren $c$.
Las personas cumplen
\[
 b_{Q_p}=0,\qquad
 \operatorname{AllDifferent}((Y_p)_p),\qquad
 \operatorname{AllDifferent}((X_p)_p).
\]
Finalmente $K_p=r_{Q_p}$, $h_p\Longleftrightarrow K_p=K_v$ y
$\sum_{p\neq v}h_p=1$: la habitaci\'on de la v\'ictima contiene
exactamente dos personas.

Es un problema de factibilidad sin objetivo, resuelto con b\'usqueda
aleatorizada. Su formulaci\'on detallada est\'a en
\texttt{generate-solution.tex}. Una vez extra\'ido $S^\star$, el mapa
y todos los muebles quedan fijos para las fases siguientes.

\section{Fase 2: construcci\'on del pool}

Sea $\mathcal R$ el pool final de pistas. Para una pista $r$:
\begin{description}
  \item[$p(r)$] Su \'unico protagonista, el sujeto \texttt{person} de
        la cl\'ausula, aunque se mencionen otras personas.
  \item[$d_r$] Su dificultad entera positiva.
  \item[$\varphi_r(S)$] Su predicado de verdad sobre una colocaci\'on.
\end{description}

Se enumeran las familias habilitadas por \texttt{rule\_types}, usando
habitaciones, elementos, coordenadas y diferencias de coordenadas
presentes en el mapa. No se generan pistas cuyo sujeto sea la
v\'ictima, pero esta s\'i puede ser una persona objetivo.

Las reglas simples se conservan si $\varphi_r(S^\star)=1$.
Si est\'a habilitada \texttt{any\_of}, se forman pares no ordenados
de condiciones simples habilitadas, sin anidamiento y con el mismo
sujeto. Solo se conserva el par si
\[
 \varphi_{r_1}(S^\star)+\varphi_{r_2}(S^\star)=1.
\]
Este filtro hace que exactamente una rama sea verdadera
\emph{en el objetivo}. La sem\'antica posterior sigue siendo
$\varphi_{r_1}\lor\varphi_{r_2}$: no se impone XOR en las alternativas.

Las dificultades implementadas son:
\begin{center}
\begin{tabular}{ll}
\hline
Familia & Dificultad $d_r$\\
\hline
\texttt{in} & $1$; $2$ con exclusividad\\
\texttt{in\_coordinate} & $2$\\
\texttt{next\_to} & $2$\\
\texttt{not\_next\_to} & $3$\\
\texttt{room\_relation} & $2$ misma habitaci\'on; $3$ diferente\\
\texttt{relative\_position} &
 $3+[\text{distancia exacta}]+[\text{objetivo elemento}]$\\
\texttt{any\_of} & $1+\sum_a d_{r_a}$\\
\hline
\end{tabular}
\end{center}
Los corchetes representan indicadores booleanos. La dificultad total
es una suma de pesos de pistas, no una medida aprendida de dificultad
humana ni del tiempo de resoluci\'on.

\section{Fase 3: modelo maestro}

\subsection{Variables, expresiones y cobertura}

Sean $m=|\mathcal R|$ y $T\in\mathbb Z_{\geq0}$ el
\texttt{target\_difficulty}, obligatorio para este caso de uso.
Se crean
\[
 x_r\in\B\quad(r\in\mathcal R),\qquad
 D(x)=\sum_{r\in\mathcal R}d_rx_r,\qquad
 C(x)=\sum_{r\in\mathcal R}x_r.
\]
$x_r=1$ selecciona la pista; $D$ y $C$ son expresiones lineales.
Cada persona distinta de la v\'ictima debe ser sujeto de al menos
una pista:
\[
 \sum_{\substack{r\in\mathcal R\\p(r)=p}}x_r\geq1
 \qquad\forall p\in P\setminus\{v\}.
\]
Una menci\'on como objetivo no satisface cobertura. Esto se aplica
tambi\'en a las disyunciones. Una suma vac\'ia hace el maestro
infactible. La excepci\'on de cobertura es \texttt{victim}; no hay
un asesino independiente en este modelo.

\subsection{Goal Programming lexicogr\'afico}

Sea $D_{\max}=\sum_r d_r$. Se a\~nade una variable entera
\[
 e\in\{0,\ldots,E_{\max}\},\qquad
 E_{\max}=\max\{T,|D_{\max}-T|\},\qquad
 e=|D(x)-T|.
\]
La igualdad absoluta se implementa con \texttt{add\_abs\_equality}.
El objetivo es
\[
 \boxed{\min\ (m+1)e+C(x).}
\]
Como $0\leq C(x)\leq m$ y $e$ es entero, reducir el error en una unidad
mejora el objetivo m\'as que cualquier cambio posible en el n\'umero
de pistas. Por tanto, equivale exactamente a minimizar
lexicogr\'aficamente $(e,C(x))$. No se usa un peso fijo como $100$.

El maestro inicial solo impone cobertura y objetivo. No contiene
todas las restricciones de unicidad desde el principio. En cada
iteraci\'on se exige que CP-SAT devuelva \texttt{OPTIMAL}, no basta
con encontrar una selecci\'on factible.

\section{Subproblema de unicidad con mapa fijo}

Para la selecci\'on $x^{(k)}$, sea
$\mathcal R_k=\{r:x_r^{(k)}=1\}$.
Con $U$ las celdas ocupables, se construyen booleanos $q_{pc}$:
\begin{align}
 \sum_{c\in U}q_{pc}&=1 &&\forall p\in P,\\
 \sum_{\substack{p\in P,\ c\in U\\i(c)=i}}q_{pc}&\leq1
      &&\forall i,\\
 \sum_{\substack{p\in P,\ c\in U\\j(c)=j}}q_{pc}&\leq1
      &&\forall j,\\
 \varphi_r(q)&=1 &&\forall r\in\mathcal R_k.
\end{align}
Cada $\varphi_r$ se inyecta mediante la estrategia de la pista, con
literales auxiliares de equivalencia cuando es necesario.
Las f\'ormulas por familia est\'an en \texttt{verify-puzzle.tex}.
Las reglas espaciales conservan la misma sem\'antica que
\texttt{evaluate}, incluida la existencia de otra persona cuando
una posici\'on relativa tiene por objetivo un elemento.

Si $c_p^\star$ es la celda de $p$ en $S^\star$, se excluye la
colocaci\'on completa:
\[
 \bigvee_{p\in P}\neg q_{p,c_p^\star}
 \quad\Longleftrightarrow\quad
 \sum_{p\in P}q_{p,c_p^\star}\leq |P|-1.
\]
Se busca al menos una alternativa, sin objetivo de optimizaci\'on.
Antes de excluirlo, el checker valida que $S^\star$ es una
colocaci\'on factible del modelo recibido.

\paragraph{Alcance de la unicidad.}
Este subproblema no mueve habitaciones ni muebles y, a diferencia de
la fase espacial, no incorpora autom\'aticamente la condici\'on de
los dos ocupantes de la habitaci\'on de la v\'ictima. La prueba se
realiza con el modelo base del checker y las pistas seleccionadas.
Las alternativas pueden incumplir aquella condici\'on del generador;
no se descartan por ese motivo.

\section{Cortes l\'ogicos de Benders}

Si se obtiene una alternativa $S_{\mathrm{alt}}^{(k)}$, se eval\'ua
todo el pool y se define
\[
 B_k=\{r\in\mathcal R:\varphi_r(S_{\mathrm{alt}}^{(k)})=0\}.
\]
Se a\~nade al maestro la cl\'ausula
\[
 \boxed{\sum_{r\in B_k}x_r\geq1,}
 \qquad\text{equivalente a }\bigvee_{r\in B_k}x_r.
\]
El corte exige seleccionar alguna pista que rompa esa alternativa.
No obliga a conservar todas las pistas de la selecci\'on anterior
ni a aumentar mon\'otonamente su n\'umero.

Si $B_k=\varnothing$, incluso el pool completo admite la alternativa:
no puede garantizar unicidad y la app produce \texttt{ValueError}.
No se a\~nade un corte vac\'io para simular una soluci\'on.

\section{Algoritmo y garant\'ia de optimalidad}

\begin{enumerate}
  \item Validar $T$, la v\'ictima y la verdad de todas las pistas del
        pool en $S^\star$. Construir el maestro inicial.
  \item Resolver el maestro hasta demostrar optimalidad y extraer
        $\mathcal R_k$.
  \item Construir el checker con $\mathcal R_k$, excluir $S^\star$
        y buscar una alternativa.
  \item Si no existe alternativa, devolver la selecci\'on y terminar.
  \item Si existe y $B_k\neq\varnothing$, a\~nadir el corte y volver
        al paso 2. Si $B_k=\varnothing$, comunicar pool insuficiente.
\end{enumerate}

Cada selecci\'on que hace \'unico a $S^\star$ tiene que satisfacer
todos los cortes: si no eligiera ninguna pista de $B_k$, admitir\'ia
la alternativa que origin\'o ese corte. El maestro es, por tanto,
una relajaci\'on del problema de selecci\'on \'unica. Cuando su
\'optimo supera el checker, tambi\'en es \'optimo global entre los
subconjuntos \'unicos del pool para $(e,C)$.

La terminaci\'on matem\'atica se apoya en que el n\'umero de
colocaciones sobre un mapa finito es finito y cada corte elimina la
alternativa actual de las futuras selecciones. Esto no proporciona
un l\'imite pr\'actico peque\~no para el tiempo ni las iteraciones.
No hay l\'imite de tiempo o iteraciones configurado en esta app.

\section{Estados, progreso y resultado}

\begin{description}
  \item[Maestro infactible:] Se comunica \texttt{ValueError}; el pool
        no permite la selecci\'on exigida. En el modelo inicial puede
        deberse a falta de cobertura.
  \item[Maestro sin \'optimo demostrado:] Se produce
        \texttt{RuntimeError}, incluso si devuelve \texttt{FEASIBLE}.
  \item[Subproblema sin alternativa:] Tras una resoluci\'on concluyente,
        se garantiza la unicidad de $S^\star$.
  \item[Subproblema no concluyente:] Se comunica un error; no se
        interpreta como ausencia de alternativas.
\end{description}

El resultado \texttt{BuiltProblem} contiene el puzzle seleccionado,
\texttt{total\_difficulty}, \texttt{difficulty\_error} e
\texttt{iterations}. El error puede ser positivo: el target no es una
restricci\'on dura y puede ser imposible alcanzarlo con unicidad.

El callback \texttt{on\_progress} recibe un \texttt{BuildProgress}
\textbf{despu\'es} de cada comprobaci\'on de unicidad. Informa del
n\'umero de candidato, dificultad, target, error absoluto, cantidad
de pistas y unicidad. No informa mientras CP-SAT sigue resolviendo
un maestro ni cuenta sus soluciones internas. Al terminar con \'exito,
si hubo $k$ candidatos comprobados, se a\~nadieron $k-1$ cortes.

El comando conserva el pool, el puzzle y las m\'etricas en
\texttt{result.json}. Tambi\'en genera \texttt{solution.html}, con
posiciones y pistas, y \texttt{puzzle.html}, con el mismo mapa y
pistas pero sin revelar las posiciones de las personas. La
presentaci\'on no construye modelos matem\'aticos adicionales.

\end{document}
